\section{A targeted response to the minibatch confound}
\label{sec:fullbatch}

We did not stop at documenting the different tail batches in the original
replay comparison. A separately frozen follow-up clones the primary runner
and changes the fitting batch assembly: carry the end of a shuffled pool
into a new independent permutation until a full 128-example batch is formed.
Only the final arithmetic-budget remainder can form a smaller batch. All
models, task laws, proposal rules, optimizer settings, data arrivals and
resource charges remain unchanged. The original source and outcomes are
preserved. This is a targeted post-hoc investigation, not a retroactive
replacement of the main experiment.

Before running seeds 850000--850011, we froze the derived source, four-arm
configuration and protocol at 05:49:36 UTC on 20 September 2026. The arms are
historical moments with replay 128, random growth with replay 128 and 138,
and fixed width seven with replay 138. We omit the separate paired-branch
diagnostic here, as specified before execution. The 48 new configurations
produce 192 complete trajectories and 1,344 events in 9.76 seconds. Every
random-128 and random-138 run now has exactly 975 optimizer updates and
124,285 fitting-example accesses. Their example counts, update counts and
minibatch-size sequences are thus matched, while their retained samples can
still differ. Historical moments spend candidate work and use 892 updates;
fixed width seven uses 392 updates, reflecting its larger per-example cost.

\begin{table}[tb]
\centering\small
\begin{tabular}{lrrrr}
\toprule
Method & Null & Rank two & Indefinite & Nonlinear \\
\midrule
Historical moment & 0.260777 & 0.000022 & 0.096205 & 0.162987 \\
Random R128 & 0.259258 & 0.000047 & 0.098319 & 0.161161 \\
Random R138 & 0.259687 & 0.000012 & 0.096821 & 0.163098 \\
Fixed width 7 & 0.259182 & 0.000021 & 0.104035 & 0.159968 \\
\bottomrule
\end{tabular}

\caption{New-seed, full-minibatch follow-up. These are separate observations,
not overwritten entries in Table~\ref{tab:quadratic-results}. All means are
final population MSE, with the same nonlinear Monte Carlo evaluation rule.}
\label{tab:fullbatch}
\end{table}

The original adverse indefinite-task comparison between replay sizes does
not persist. Random-128 minus random-138 is now $+.001498$, with interval
$[-.001121,.004117]$, compared with the original negative estimate
$-.005038$ and interval $[-.007834,-.002243]$. Historical moments minus
random-138 is $-.000616$, with interval $[-.001736,.000504]$; historical
moments minus random-128 is $-.002114$, with interval
$[-.004453,.000226]$. Both intervals include zero.
Thus a general learning benefit from retaining the target moment is not
established by this follow-up. On the other tasks, all targeted historical
and replay-size paired intervals likewise include zero. That statement is a
failure to resolve effects, not an equivalence proof.

The changed update schedule is mechanically verified, and the new empirical
comparisons fail to support the original strong attribution. Because the
follow-up uses new seeds, the difference between old and new point estimates
is not itself a paired causal estimate of the minibatch correction. The
correct conclusion is narrower: the original comparison was confounded,
and the proposed moment advantage has not survived the present control with
matched minibatch schedules. A stronger assessment would randomize or pair
both batching schedules on an additional untouched seed set, ideally with
an optimizer-budget definition chosen in advance.

\begin{figure}[tb]
\centering\includegraphics[width=\linewidth]{fullbatch_followup.pdf}
\caption{Targeted paired risk differences after removing epoch-tail changes
in update count. Bars are descriptive 95\% Student-$t$ intervals over twelve
independent new seeds; no multiplicity adjustment is made. None of these
intervals establishes a resolved final-risk advantage.}
\end{figure}

Across the primary, adaptive and targeted studies there are 768 complete
confirmation/follow-up trajectories, plus 1,440 same-checkpoint trained
branches. This volume is evidence of implementation and investigation,
not a substitute for a surviving theorem-level novelty claim or a robust
empirical improvement. The unsuccessful attribution remains part of the
research result.
